\documentclass{article}
\usepackage[T1]{fontenc}
\usepackage{lmodern}
\usepackage{graphicx}
\usepackage{amsmath,amssymb}
\usepackage[a4paper,margin=2cm]{geometry}
\usepackage[absolute,overlay]{textpos}
\setlength{\TPHorizModule}{1mm}
\setlength{\TPVertModule}{1mm}
\usepackage[indonesian]{babel}
\usepackage{hyperref}
\setlength{\emergencystretch}{2em}
% unit: Halaman 13

\begin{document}

% === Halaman 13 (python/native) ===

13

• Contoh:

          Plainteks: 10100010001110101001

     Bagi plainteks menjadi blok-blok 4-bit:


1010  0010  0011  1010  1001

     ( dalam notasi HEX :A23A9)

• Kunci (juga 4-bit): 1011 

• Misalkan fungsi enkripsi E yang sederhana  algoritmanya sbb:

   1.  XOR-kan blok plainteks Pi dengan K 

   2.  geser secara wrapping bit-bit dari hasil langkah 1 satu posisi ke kiri


EK(P) = (P  K) << 1

\end{document}
